\documentclass[10pt,twocolumn]{article}

\usepackage[margin=0.75in]{geometry}
\usepackage{graphicx}
\usepackage{booktabs}
\usepackage{amsmath}
\usepackage{hyperref}
\usepackage{caption}
\usepackage{subcaption}
\usepackage{float}
\usepackage{enumitem}
\usepackage{parskip}
\usepackage[numbers]{natbib}

\setlength{\columnsep}{0.25in}
\captionsetup{font=small, labelfont=bf}

% ── Figure root ───────────────────────────────────────────────────────────────
% For Overleaf: figures sit at zip root alongside report.tex
% For local build from docs/final_report/: use {../../outputs/} instead
\graphicspath{{./}}

\title{\textbf{Digit Detection and Recognition Using\\
Convolutional Neural Networks}\\
\large GT CS 6476 --- Computer Vision Final Project}

\author{Siddhanth Satish}
\date{}

\begin{document}
\maketitle

% ─────────────────────────────────────────────────────────────────────────────
\section{Introduction}
% ─────────────────────────────────────────────────────────────────────────────

Recognising sequences of digits in natural scene images is a fundamental
computer vision problem with applications in autonomous navigation, postal
automation, and document digitisation.
Unlike isolated digit classification on clean crops, scene-level digit
detection must contend with scale variation, arbitrary orientation, uneven
illumination, background clutter, and sensor noise---all simultaneously.

This project implements a complete two-stage pipeline:
(1) a region-proposal stage that generates candidate bounding boxes using
Maximally Stable Extremal Regions (MSER)~\cite{matas2004robust} and a
sliding-window image pyramid, and
(2) a classification stage that labels each proposal as one of ten digit
classes or background using a Convolutional Neural Network (CNN).
Two classifiers are evaluated: a custom CNN trained from scratch and a
fine-tuned VGG16~\cite{simonyan2015vgg} initialised with ImageNet weights.
Post-processing applies per-class Non-Maximum Suppression (NMS) and
left-to-right sorting to assemble the final digit sequence.

All models are trained on the Street View House Numbers (SVHN) Format~1
dataset~\cite{netzer2011svhn}, which provides 73,257 training and 26,032
test digit crops extracted from Google Street View imagery.

% ─────────────────────────────────────────────────────────────────────────────
\section{Related Work}
% ─────────────────────────────────────────────────────────────────────────────

\paragraph{CNN-based digit recognition.}
LeCun et al.~\cite{lecun1998gradient} introduced LeNet, demonstrating that
convolutional architectures could learn spatial hierarchies of features
directly from pixel data.
Goodfellow et al.~\cite{goodfellow2014multidigit} extended this to
multi-digit sequence recognition in a single forward pass, achieving
near-human accuracy on SVHN.
More recently, He et al.~\cite{he2016resnet} showed that residual
connections allow very deep networks to be trained stably, further pushing
classification accuracy.

\paragraph{Transfer learning.}
Simonyan and Zisserman~\cite{simonyan2015vgg} demonstrated that deep
uniform architectures (VGG) pre-trained on ImageNet transfer effectively
to new visual tasks.
Yosinski et al.~\cite{yosinski2014transferable} provided theoretical and
empirical evidence that lower convolutional layers learn general-purpose
edge and texture detectors reusable across domains.

\paragraph{Region proposals.}
Matas et al.~\cite{matas2004robust} introduced MSER as a method for
detecting stable connected components across a range of intensity
thresholds, making it well-suited for text-like regions.
Viola and Jones~\cite{viola2004robust} popularised the sliding-window
paradigm with image pyramids for scale-invariant detection.
Ren et al.~\cite{ren2015fasterrcnn} later unified region proposal and
classification in Faster R-CNN; however, the assignment prohibits
one-stage detectors, motivating the classical two-stage approach used here.

% ─────────────────────────────────────────────────────────────────────────────
\section{Method}
% ─────────────────────────────────────────────────────────────────────────────

\subsection{Preprocessing}

Every input image first undergoes \textbf{polarity normalisation}: if the
median pixel brightness exceeds 127 the image is inverted, ensuring that
digit strokes are always brighter than the background.
This stabilises MSER detection, which relies on consistent contrast polarity.

The image is then denoised with a $3\times3$ Gaussian blur to suppress
high-frequency sensor noise.
Lighting normalisation is applied via Contrast Limited Adaptive Histogram
Equalisation (CLAHE)~\cite{zuiderveld1994clahe} on the L channel of the
LAB colour space (\texttt{clipLimit}$=2.0$,
\texttt{tileGridSize}$=8\times8$).
Operating in LAB isolates luminance from chrominance, preventing colour
distortion while improving local contrast.

\begin{figure}[H]
  \centering
  \includegraphics[width=\linewidth]{report_figures/comparison/pipeline_steps.png}
  \caption{Full preprocessing pipeline for one correctly detected image
           (top) and one failure case (bottom).
           Steps: original $\to$ polarity normalisation $\to$ Gaussian blur
           $\to$ CLAHE $\to$ MSER proposals $\to$ deduplication $\to$
           classifier hits $\to$ DBSCAN clustering $\to$ final detections.}
  \label{fig:pipeline}
\end{figure}

\subsection{Region Proposals}

\paragraph{MSER.}
MSER detects connected pixel regions that remain stable across a wide range
of intensity thresholds.
Digit strokes exhibit this property because they maintain consistent
contrast against their background.
MSER is run on both the greyscale image and its inverse to capture both
dark-on-light and light-on-dark digits.
Regions with area outside $[50, 3000]$ pixels or aspect ratio outside
$[0.2, 5.0]$ are discarded as unlikely to contain a single digit.

\paragraph{Sliding window + image pyramid.}
A $32\times32$ window slides across the image in steps of 24 pixels.
The image is repeatedly downscaled by a factor of 0.9 until either
dimension falls below 32 pixels, producing a multi-scale pyramid.
All proposals are mapped back to original-image coordinates by dividing by
the cumulative scale factor.
By default, only MSER proposals are used at inference
(\texttt{use\_sliding}$=$\texttt{false} in \texttt{config.json}); the
sliding window can be enabled for higher recall at the cost of
significantly more proposals.
Near-duplicate proposals (IoU $> 0.8$) are removed before classification.

\subsection{Classification}

\paragraph{11-class formulation.}
Both models output 11 classes: digits 0--9 (indices 0--9) and a background
class (index 10).
The background class is trained on synthetic crops: between-digit patches
sampled from SVHN scene images (50\%), random noise patches (25\%), and
linear gradient patches (25\%).
The ratio of background to digit samples is controlled by
\texttt{bg\_ratio}$=1.0$ in \texttt{config.json}, producing a balanced
dataset with one background crop per digit crop.
At inference, any proposal classified as background is discarded before NMS.

\paragraph{CustomCNN.}
The custom architecture consists of three convolutional blocks, each
containing two $3\times3$ Conv--BatchNorm--ReLU layers followed by
$2\times2$ MaxPool and Dropout2d(0.25).
Filter counts double across blocks: $32\to64\to128$.
The classifier head is Linear($2048\to512$)--ReLU--Dropout(0.5)--Linear($512\to11$).
The network is trained from random initialisation using Adam
($\text{lr}=10^{-3}$, weight decay $10^{-4}$, batch size 128) with
ReduceLROnPlateau (patience 3, factor 0.5) and early stopping (patience 7).

\paragraph{VGG16 fine-tuning.}
VGG16~\cite{simonyan2015vgg} is loaded with ImageNet weights.
The final fully-connected layer is replaced with Linear($4096\to11$).
All layers are unfrozen; the convolutional backbone uses a learning rate
ten times smaller than the classifier head ($10^{-5}$ vs.\ $10^{-4}$) to
preserve pretrained features while allowing domain adaptation.
Batch size is 64 to accommodate GPU memory.

\paragraph{Data augmentation.}
Training crops are augmented with ColorJitter (brightness/contrast $\pm0.3$,
saturation $\pm0.2$) and RandomAffine (rotation $\pm15^\circ$, translation
$\pm10\%$, scale $[0.9, 1.1]$) to improve robustness to the pose and
lighting variations expected at test time.
Validation and test sets use no augmentation.

\subsection{Post-processing}

After classification, proposals with confidence below 0.85 or predicted as
background are removed.
Per-class NMS with IoU threshold 0.3 suppresses duplicate detections of the
same digit.
Running NMS per class prevents a high-confidence ``4'' from suppressing a
nearby ``2''.
Surviving detections are sorted by their left-edge $x$-coordinate to
assemble the digit sequence left-to-right.

% ─────────────────────────────────────────────────────────────────────────────
\section{Experiments}
% ─────────────────────────────────────────────────────────────────────────────

\subsection{Dataset and Evaluation}

SVHN Format~1 provides full street-view images with per-digit bounding
boxes stored in \texttt{digitStruct.mat}.
Each bounding box is cropped and resized to $32\times32$.
The training set (73,257 crops) is split 95/5 into train and validation
subsets; the test set contains 26,032 digit crops.
With \texttt{bg\_ratio}$=1.0$, an equal number of synthetic background
crops are added, yielding 52,064 total test samples (26,032 digits +
26,032 background).
Pixel values are normalised using SVHN channel statistics
($\mu=[0.4377, 0.4438, 0.4728]$, $\sigma=[0.198, 0.201, 0.197]$).

The SVHN test set exhibits significant class imbalance: digit~1 is the
most frequent class ($\approx$3$\times$ more samples than digit~9), and
the majority of images contain 2-digit sequences (8,356 of 13,068 images).

\begin{figure}[H]
  \centering
  \begin{subfigure}[b]{0.48\linewidth}
    \includegraphics[width=\linewidth]{report_figures/exploratory/eda_class_distribution.png}
    \caption{Per-digit class distribution}
  \end{subfigure}
  \hfill
  \begin{subfigure}[b]{0.48\linewidth}
    \includegraphics[width=\linewidth]{report_figures/exploratory/eda_sequence_length.png}
    \caption{Sequence length distribution}
  \end{subfigure}
  \caption{SVHN dataset statistics. Digit~1 dominates both splits;
           2-digit sequences are most common.}
  \label{fig:eda}
\end{figure}

\subsection{Training Curves}

\begin{figure}[H]
  \centering
  \includegraphics[width=\linewidth]{report_figures/comparison/comparison_loss.png}
  \caption{Train and validation cross-entropy loss for both models.
           VGG16 converges faster due to pretrained initialisation.}
  \label{fig:loss}
\end{figure}

\begin{figure}[H]
  \centering
  \includegraphics[width=\linewidth]{report_figures/comparison/comparison_accuracy.png}
  \caption{Train and validation accuracy. CustomCNN requires more epochs
           to reach comparable accuracy.}
  \label{fig:acc}
\end{figure}

Figure~\ref{fig:loss} shows that VGG16 achieves low validation loss by
epoch~5, whereas CustomCNN requires approximately 20 epochs to reach a
similar regime.
The learning rate schedule for CustomCNN decays from $10^{-3}$ to
$5\times10^{-4}$ after epoch~24 (one plateau reduction), while VGG16
decays from $10^{-4}$ to $5\times10^{-5}$ after epoch~14.
Both models show a small but consistent generalisation gap (train accuracy
exceeds validation accuracy by $\approx$0.5--1.5 percentage points),
indicating mild overfitting that is well-controlled by dropout and weight
decay.

\subsection{Crop-Level Classification Results}

\begin{table}[H]
  \centering
  \caption{Crop-level test set performance (52,064 samples: 26,032 digit
           crops + 26,032 synthetic background crops).}
  \label{tab:results}
  \begin{tabular}{lccc}
    \toprule
    Model & Train Acc & Val Acc & Test Acc \\
    \midrule
    CustomCNN   & 97.99\% & 98.45\% & 96.06\% \\
    VGG16       & 98.93\% & 98.69\% & \textbf{97.02\%} \\
    \bottomrule
  \end{tabular}
\end{table}

Table~\ref{tab:results} shows that VGG16 outperforms CustomCNN by
0.96 percentage points on the test set.
The pretrained backbone provides a stronger initialisation and converges
in fewer epochs (20 vs.\ 30), confirming the benefit of transfer learning
even for a relatively simple $32\times32$ classification task.

\begin{table}[H]
  \centering
  \caption{Per-class precision / recall / F1 on the test set.}
  \label{tab:perclass}
  \resizebox{\linewidth}{!}{%
  \begin{tabular}{l ccc ccc}
    \toprule
    & \multicolumn{3}{c}{CustomCNN} & \multicolumn{3}{c}{VGG16} \\
    \cmidrule(lr){2-4} \cmidrule(lr){5-7}
    Class & Prec & Rec & F1 & Prec & Rec & F1 \\
    \midrule
    0  & 0.88 & 0.97 & 0.92 & 0.92 & 0.97 & 0.94 \\
    1  & 0.91 & 0.94 & 0.93 & 0.92 & 0.96 & 0.94 \\
    2  & 0.97 & 0.97 & 0.97 & 0.97 & 0.97 & 0.97 \\
    3  & 0.94 & 0.95 & 0.94 & 0.96 & 0.94 & 0.95 \\
    4  & 0.91 & 0.97 & 0.94 & 0.96 & 0.98 & 0.97 \\
    5  & 0.94 & 0.96 & 0.95 & 0.95 & 0.97 & 0.96 \\
    6  & 0.94 & 0.96 & 0.95 & 0.95 & 0.97 & 0.96 \\
    7  & 0.90 & 0.95 & 0.93 & 0.95 & 0.94 & 0.95 \\
    8  & 0.91 & 0.96 & 0.94 & 0.92 & 0.97 & 0.94 \\
    9  & 0.92 & 0.97 & 0.94 & 0.93 & 0.97 & 0.95 \\
    bg & 1.00 & 0.96 & 0.98 & 1.00 & 0.97 & 0.99 \\
    \bottomrule
  \end{tabular}}
\end{table}

Table~\ref{tab:perclass} shows that VGG16 improves precision across
nearly all digit classes compared to CustomCNN.
Digit~0 has the lowest precision for CustomCNN (0.88), likely because the
circular stroke of ``0'' resembles ``6'' and ``8'' in small, noisy crops.
The background class achieves precision 1.00 for both models, confirming
that the synthetic background training data is effective.

\begin{figure}[H]
  \centering
  \begin{subfigure}[b]{0.48\linewidth}
    \includegraphics[width=\linewidth]{results/vgg16_confusion_matrix.png}
    \caption{VGG16}
  \end{subfigure}
  \hfill
  \begin{subfigure}[b]{0.48\linewidth}
    \includegraphics[width=\linewidth]{results/custom_confusion_matrix.png}
    \caption{CustomCNN}
  \end{subfigure}
  \caption{Confusion matrices on the test set. Both models most frequently
           confuse visually similar digit pairs (0/6, 7/1, 3/8).}
  \label{fig:cm}
\end{figure}

\begin{figure}[H]
  \centering
  \includegraphics[width=\linewidth]{report_figures/comparison/comparison_bar.png}
  \caption{Train, validation, and test accuracy comparison.}
  \label{fig:bar}
\end{figure}

\begin{figure}[H]
  \centering
  \begin{subfigure}[b]{0.48\linewidth}
    \includegraphics[width=\linewidth]{results/custom_per_class_accuracy.png}
    \caption{CustomCNN}
  \end{subfigure}
  \hfill
  \begin{subfigure}[b]{0.48\linewidth}
    \includegraphics[width=\linewidth]{results/vgg16_per_class_accuracy.png}
    \caption{VGG16}
  \end{subfigure}
  \caption{Per-class accuracy on the test set. VGG16 is more uniform
           across classes; CustomCNN shows a larger spread.}
  \label{fig:perclass_acc}
\end{figure}

\subsection{Region Proposal Evaluation}

Before evaluating end-to-end sequence accuracy, we assess the quality of
the region proposal stage in isolation.
A ground-truth bounding box is considered ``hit'' if any proposal overlaps
it with IoU $\geq 0.3$.
An image has ``full coverage'' if every GT box in that image is hit.

\begin{table}[H]
  \centering
  \caption{Region proposal quality on the SVHN test set (13,068 images,
           26,032 GT boxes, 809,812 proposals).}
  \label{tab:proposals}
  \begin{tabular}{lc}
    \toprule
    Metric & Value \\
    \midrule
    Recall (GT boxes hit)       & 93.67\% \\
    Precision (proposals useful) & 15.07\% \\
    Full coverage (all GT hit)  & 89.56\% \\
    \bottomrule
  \end{tabular}
\end{table}

\begin{figure}[H]
  \centering
  \includegraphics[width=\linewidth]{report_figures/comparison/proposal_hit_rate.png}
  \caption{Proposal recall, precision, and full-coverage rate broken down
           by ground-truth sequence length. Recall remains above 89\% for
           all lengths, but full coverage drops for longer sequences.}
  \label{fig:proposal_hit}
\end{figure}

Table~\ref{tab:proposals} shows that MSER achieves 93.67\% recall,
meaning the vast majority of digits receive at least one overlapping
proposal.
However, precision is only 15.07\%---most proposals cover background
regions---which places a heavy burden on the classifier to reject
non-digit crops.
Full coverage drops from 97.0\% for 1-digit images to 77.4\% for 4-digit
images (Figure~\ref{fig:proposal_hit}), explaining much of the
sequence-level accuracy degradation for longer sequences.

\subsection{Sequence-Level Evaluation}

The ultimate metric for this pipeline is sequence-level accuracy: the
predicted digit string must exactly match the ground-truth string for an
image to be counted as correct.

\begin{table}[H]
  \centering
  \caption{Sequence-level accuracy on the SVHN test set (13,068 images).}
  \label{tab:seq}
  \begin{tabular}{lcccc}
    \toprule
    Model & Correct & Wrong digits & Wrong length & No detection \\
    \midrule
    CustomCNN & \textbf{44.08\%} & 9.67\% & 43.85\% & 2.39\% \\
    VGG16     & 43.14\% & 9.93\% & 44.86\% & 2.07\% \\
    \bottomrule
  \end{tabular}
\end{table}

\begin{figure}[H]
  \centering
  \includegraphics[width=\linewidth]{report_figures/comparison/comparison_seq_accuracy.png}
  \caption{Sequence-level accuracy and error breakdown. ``Wrong length''
           dominates errors for both models.}
  \label{fig:seq_acc}
\end{figure}

\begin{figure}[H]
  \centering
  \includegraphics[width=\linewidth]{report_figures/comparison/comparison_seq_by_length.png}
  \caption{Sequence accuracy vs.\ ground-truth sequence length. Accuracy
           drops sharply for sequences of 3+ digits.}
  \label{fig:seq_len}
\end{figure}

Table~\ref{tab:seq} reveals that wrong-length predictions account for
$\approx$44\% of all test images---the dominant failure mode.
This occurs when the proposal stage either misses a digit entirely (under-
detection) or generates spurious detections that survive NMS (over-detection).
Only $\approx$10\% of errors are wrong-digit errors (correct length but
incorrect digit identity), confirming that the classifier is accurate when
it receives a good crop.

Interestingly, CustomCNN achieves slightly higher sequence accuracy
(44.08\%) than VGG16 (43.14\%) despite lower crop-level test accuracy
(96.06\% vs.\ 97.02\%).
This suggests that CustomCNN's confidence calibration interacts more
favourably with the fixed confidence threshold (0.85) and NMS pipeline,
producing fewer spurious detections.

Figure~\ref{fig:seq_len} shows that accuracy degrades sharply with
sequence length: 62.4\% for 1-digit sequences, 44.6\% for 2-digit, 23.1\%
for 3-digit, and only 3.4\% for 4-digit sequences.
This is consistent with the proposal full-coverage drop-off
(Figure~\ref{fig:proposal_hit}): missing even one digit in a multi-digit
sequence causes a wrong-length error.

\subsection{Detection Results}

\begin{figure}[H]
  \centering
  \begin{subfigure}[b]{0.48\linewidth}
    \includegraphics[width=\linewidth]{graded_images/vgg16/1.png}
    \caption{Scale invariance}
  \end{subfigure}
  \hfill
  \begin{subfigure}[b]{0.48\linewidth}
    
\includegraphics[width=\linewidth]{graded_images/vgg16/2.png}
    \caption{Orientation ($30^\circ$)}
  \end{subfigure}
  \vspace{4pt}
  \begin{subfigure}[b]{0.48\linewidth}
    \includegraphics[width=\linewidth]{graded_images/vgg16/3.png}
    \caption{Corner location}
  \end{subfigure}
  \hfill
  \begin{subfigure}[b]{0.48\linewidth}
    
\includegraphics[width=\linewidth]{graded_images/vgg16/4.png}
    \caption{Low lighting}
  \end{subfigure}
  \vspace{4pt}
  \begin{subfigure}[b]{0.48\linewidth}
    
\includegraphics[width=\linewidth]{graded_images/vgg16/5.png}
    \caption{Gaussian noise}
  \end{subfigure}
  \caption{VGG16 detection results on the five graded demo images.
           Green boxes show accepted detections with confidence scores.}
  \label{fig:detections}
\end{figure}

\paragraph{Positive results.}
The pipeline correctly detects digits at multiple scales (Figure~\ref{fig:detections}a)
because the image pyramid ensures that small digits become full-sized
proposals at lower pyramid levels.
CLAHE normalisation enables detection under low-light conditions
(Figure~\ref{fig:detections}d) by adaptively boosting local contrast before
region proposal.
The Gaussian blur pre-processing step suppresses noise artefacts
(Figure~\ref{fig:detections}e) that would otherwise generate spurious MSER
regions.

\paragraph{Negative results and failure analysis.}
Rotated digits (Figure~\ref{fig:detections}b) present the most significant
challenge.
The sliding window always proposes axis-aligned boxes, so a digit rotated
by $30^\circ$ may be cropped with substantial background padding, reducing
the effective digit area in the $32\times32$ input.
Training augmentation covers $\pm15^\circ$, but larger rotations fall
outside the training distribution.

Digits placed in extreme corners (Figure~\ref{fig:detections}c) can be
partially missed if the sliding window step size causes the window to
straddle the image boundary, producing a crop that is mostly background.

The sequence-level results (Section~4.5) confirm that the dominant failure
mode is wrong-length prediction ($\approx$44\% of images), driven by the
proposal stage rather than the classifier.
Per-class NMS with IoU threshold 0.3 may suppress one detection when two
digits are very close together and their bounding boxes overlap
significantly.
This is an inherent limitation of the two-stage approach; end-to-end
detectors such as Faster R-CNN~\cite{ren2015fasterrcnn} handle this more
gracefully through anchor regression.

% ─────────────────────────────────────────────────────────────────────────────
\section{Conclusion}
% ─────────────────────────────────────────────────────────────────────────────

This project implemented a complete digit detection and recognition system
combining classical region proposal (MSER + sliding-window pyramid) with
deep CNN classification.
Two models were trained on SVHN: a custom CNN from scratch achieving
96.06\% crop-level test accuracy, and a fine-tuned VGG16 achieving 97.02\%.
The 11-class formulation with a synthetic background class enables explicit
rejection of non-digit proposals, reducing false positives compared to
confidence-threshold-only filtering.

VGG16 converges faster and achieves higher crop-level accuracy, confirming
the benefit of transfer learning from ImageNet.
However, at the sequence level both models achieve similar accuracy
($\approx$44\%), with CustomCNN marginally outperforming VGG16 (44.08\%
vs.\ 43.14\%).
The dominant failure mode is wrong-length prediction caused by missed or
spurious region proposals, not misclassification of individual digits.

The primary bottleneck is the region proposal stage: while MSER achieves
93.67\% per-box recall, full-coverage drops to 77\% for 4-digit sequences,
directly limiting sequence accuracy.
Future work could replace the hand-crafted proposal stage with a learned
region proposal network, incorporate rotation-invariant features, or adopt
an end-to-end detector to address these limitations.

% ─────────────────────────────────────────────────────────────────────────────
\bibliographystyle{plainnat}
\bibliography{references}

\end{document}
